\begin{abstract}
\noindent
Tropical peatland wildfires in Southeast Asia represent a critical land governance dilemma. While environmental enforcement emphasizes corporate concession liability, the spatial distribution of fires across jurisdictional boundaries remains contested. Here, we implement a Spatial Boundary Discontinuity Design (Spatial BDD) using 43,937 verified satellite active fire detections (NASA FIRMS) during the October 2026 crisis in Central Kalimantan, integrated with 812 audited cadastral concessions (627 oil palm HGU and 185 forestry PBPH), metric canal vectors (UTM Zone 49S), and ERA5-Land covariates. Point-in-polygon overlay reveals that 75.4\% (33,138 hotspots) ignited outside corporate perimeters on unallocated public lands (\textit{Area Penggunaan Lain}, APL), exposing an acute governance void. Within a 3-km frontier corridor, Spatial BDD establishes a sharp +56.5\% discontinuity in fire density, jumping from 0.183 hotspots/km$^2$ inside concessions to 0.286 hotspots/km$^2$ on exterior APL peatlands (peaking at 0.370 hotspots/km$^2$ at 3\,km). Under 5-fold Spatial Block Cross-Validation, coordinate-dependent machine learning collapsed due to spatial autocorrelation memorization (ROC-AUC plunging from 0.7956 to 0.5399), whereas a pure biophysical Random Forest generalized across regencies (holdout ROC-AUC: 0.7527; 2.4-fold higher spatial $F_1$-score). TreeSHAP identified a 3-km canal corridor and 14-day drought trigger as key tipping points. Corporate strict liability (Law No.~32/2009) deters burning inside concessions but leaves adjacent state peatlands vulnerable to drainage spillover. Sustainable land governance requires extending statutory buffer liability across 3\,km and formalizing community peatland management units.

\vspace{0.4em}
\noindent\textbf{Keywords:}\ \ \textit{Applied spatial geography; Tropical peatlands; Spatial boundary discontinuity; Spatial block cross-validation; Land tenure governance; Wildfire intensity modeling; Spatial autocorrelation.}
\end{abstract}
